\documentclass[11pt]{article}
\usepackage[utf8]{inputenc}
\usepackage{amsmath,amssymb}
\usepackage[margin=1in]{geometry}
\usepackage{hyperref}
\emergencystretch=3em

\title{Divisible perturbations of the weight:\\
$Z(a+u^k\tilde J,\ \eta_0+u^k\tilde\eta)\sim\eta_0\,C\,n^{-p/2}(\log n)^m$}
\author{Claude, with Daniel Murfet}
\date{2026-09-06}

\begin{document}
\maketitle
\sloppy

\begin{abstract}
The mechanism of unit~65 extends to the weight. For $\eta=\eta_0+u^k\tilde\eta$ and
$\xi=a+u^k\tilde J$ with $\tilde\eta,\tilde J$ continuous on the box, the weight perturbation carries
the extra factor $u^k=t/\sqrt n$, and $te^{-(\beta/4)t^2}\le4/\beta+1$ gives
$|Z(\xi,\eta)-\eta_0Z(\xi,1)|\le\frac{L_\eta(4/\beta+1)}{\sqrt n}Z_{\beta/2,2a}(n)$ for $n\ge16L^2$.
Combined with unit~65, $Z(\xi,\eta)-\eta_0Z(a,1)=O(n^{-1/2}Z_{\beta/2,2a})$, so for $\eta_0\ne0$ the
standard integral is $\sim\eta_0Cn^{-p/2}(\log n)^m$ with the constant-$\xi$ exponent, multiplicity and
constant: the leading coefficient is $\eta(0)$ times the population constant, as the paper's remark on
the sign of $\eta$ anticipates. Zero \texttt{sorry}/\texttt{axiom}.
\end{abstract}

\section{Setting}
Both perturbations vanish on the exceptional divisor to order $u^k$, which is exactly what makes them
lower order: after the substitution $t=\sqrt n\,u^k$ each contributes a factor $t/\sqrt n$ or
$t^2/\sqrt n$, absorbed by comparing with the kernel at the halved temperature.

\section{The things formalised}
{\small
\begin{verbatim}
noncomputable def boxIntegralXiEta (β a b c : ℝ) {d : ℕ} (k h : Fin d → ℕ)
    (J η : (Fin d → ℝ) → ℝ) : ℝ := ∫ u, (∏ i, u i ^ h i) * η u
    * Real.exp (-β * (c * ∏ i, u i ^ k i) ^ 2 + β * (c * ∏ i, u i ^ k i) * (a + J u))
    ∂(boxMeasure b d)
theorem mul_exp_neg_quarter_le (β t : ℝ) (hβ : 0 < β) :
    t * Real.exp (-(β / 4) * t ^ 2) ≤ 4 / β + 1
theorem boxIntegralXiEta_sub_le ... (hnL : 16 * L ^ 2 ≤ n) :
    |boxIntegralXiEta β a b (Real.sqrt n) k h (fun u => (∏ i, u i ^ k i) * J u)
          (fun u => η₀ + (∏ i, u i ^ k i) * η u)
        - η₀ * boxIntegralXi β a b (Real.sqrt n) k h (fun u => (∏ i, u i ^ k i) * J u)|
      ≤ Lη * (4 / β + 1) / Real.sqrt n * boxIntegralFin (β / 2) (2 * a) b (Real.sqrt n) k h
theorem boxIntegralXiEta_isEquivalent ... (hη₀ : η₀ ≠ 0) :
    ∃ (p : ℝ) (m : ℕ) (C : ℝ), 0 < C ∧ (∀ i, p ≤ finExp k h i) ∧
      m + 1 = (Finset.univ.filter fun i => finExp k h i = p).card ∧
      (fun n => boxIntegralXiEta β a b (Real.sqrt n) k h (fun u => (∏ i, u i ^ k i) * J u)
          (fun u => η₀ + (∏ i, u i ^ k i) * η u))
        ~[atTop] fun n => η₀ * C * (n ^ (-(p / 2)) * Real.log n ^ m)
\end{verbatim}
}

\section{Proof strategy}
Pointwise, the weight-perturbation integrand is $u^h\,(t/\sqrt n)\tilde\eta\,e^{-\beta t^2+\beta t(a+t\tilde J/\sqrt n)}$;
the exponential is at most $e^{-(\beta/4)t^2}e^{-(\beta/2)t^2+\beta ta}$ when $L/\sqrt n\le1/4$, and
$te^{-(\beta/4)t^2}\le4/\beta+1$ follows from $t\le t^2+1$ and $t^2\le\frac4\beta e^{(\beta/4)t^2}$
(\texttt{Real.add\_one\_le\_exp}). Integration as in unit~65. For the asymptotic the two differences
are combined by the triangle inequality into an $O(n^{-1/2}Z_{\beta/2,2a})$ bound, which is
$o(\eta_0Z(a))$; \texttt{IsEquivalent.add\_isLittleO} concludes after multiplying unit~59's
equivalence by the constant $\eta_0$.

\section{Roadblocks and resolutions}
\begin{itemize}
\item Rewriting $t^2$ as $\frac4\beta(\frac\beta4t^2)$ with \texttt{rw} also rewrote inside the
exponential; a \texttt{calc} step instead.
\item Dot notation on an \texttt{IsEquivalent} fact fell through to \texttt{IsLittleO.const\_mul\_left}
(the definition unfolds); multiply by a constant via \texttt{IsEquivalent.refl.mul} and \texttt{congr}.
\item \texttt{rw [one\_div]} hit the exponent $1/2$ before the intended $1/\sqrt n$; use
\texttt{(one\_div \_).symm} as a term.
\end{itemize}

\section{What was Mathlib, what was new}
Mathlib: \texttt{Real.add\_one\_le\_exp}, \texttt{Real.one\_le\_exp}, \texttt{abs\_add\_le},
\texttt{IsEquivalent.add\_isLittleO}, \texttt{IsLittleO.mul\_isBigO}. New: the weight-perturbation bound
and the joint theorem for divisible $\xi$ and $\eta$ perturbations.

\section{Lessons}
Once one perturbation is handled by comparison with a halved-temperature kernel, further perturbations
of the same divisibility type follow the same template; the triangle inequality assembles them.

\section{Outcome}
The multiplicity theorem now covers weights and fluctuation functions of the form
$\eta_0+u^k\tilde\eta$, $a+u^k\tilde J$, with leading coefficient $\eta_0$ times the population
constant. What remains for general analytic $\xi,\eta$ is the genuine leading-coefficient dependence
on their restrictions to the minimal divisor components.
\end{document}
